
\documentclass[8pt,a4paper]{article}

% ------------------------------------------------
% Packages
% ------------------------------------------------
\usepackage[
    top=0.15in,
    bottom=0.15in,
    left=0.55in,
    right=0.55in
]{geometry}

\usepackage{amsmath,amsfonts,amssymb}
\usepackage{booktabs}
\usepackage{tabularx}
\usepackage{makecell}
\usepackage{caption}
\usepackage{hyperref}
\usepackage{titlesec}

% ------------------------------------------------
% Formatting
% ------------------------------------------------

% Black-and-white hyperlinks
\hypersetup{
    colorlinks=true,
    allcolors=black
}

% Reduce line spacing
\linespread{0.52}

% Reduce paragraph spacing
\setlength{\parindent}{0pt}
\setlength{\parskip}{1pt}

% Reduce section heading spacing
\titlespacing*{\subsection}
{0pt}{7pt}{3pt}

% Table formatting
\setlength{\tabcolsep}{5pt}
\renewcommand{\arraystretch}{1.05}

\captionsetup[table]{
    font=small,
    labelfont=bf,
    position=top,
    skip=3pt
}

% ------------------------------------------------
% Document
% ------------------------------------------------
\begin{document}
\thispagestyle{empty}

% ------------------------------------------------
% Title and Team Information
% ------------------------------------------------
\begin{center}

{\large\bfseries
[Project Title]}\\[4pt]

\textbf{Team No:} XX
\quad | \quad
\textbf{Members:}
Student 1 (Roll No), Student 2 (Roll No)\\[3pt]

\textbf{GitHub:}
\url{https://github.com/username/project-repo}

\end{center}

\vspace{-5pt}

% ------------------------------------------------
\subsection*{I. Abstract (Maximum 100 Words)}
% ------------------------------------------------
\begin{enumerate}
    \item State the fundamental challenge, limitation, or failure mode in existing literature that this paper targets.
    \item Explain why conventional paradigms fall short and why a new formulation or perspective is critically needed.
    \item Describe the authors' core conceptual idea and mechanism designed to overcome this bottleneck.
    \item Summarize your technical interpretation of what this approach conceptually enables beyond baseline methods.
\end{enumerate}

% ------------------------------------------------
\subsection*{II. Proposed Novelty \& Technical Contribution}
% ------------------------------------------------
\begin{enumerate}
    \item Define the exact technical novelty introduced by your team into the assigned research paper.
    \item Specify whether the novelty is an architectural modification, a new module, an altered loss function, or an algorithmic enhancement.
    \item Explicitly articulate how your contribution fundamentally differs from the baseline methodology.
\end{enumerate}

% ------------------------------------------------
\subsection*{III. Theoretical / Mathematical Justification}
% ------------------------------------------------
\begin{enumerate}
    \item Formulate the core mathematical equation, loss term, or operation governing your proposed modification.
    \item Define each variable and provide theoretical reasoning or a derivation showing why this modification improves feature representation or optimization.
    \item Explicitly establish the causal link between your mathematical formulation and the expected performance gain without making unsupported claims.
\end{enumerate}

% ------------------------------------------------
\subsection*{IV. Comparative Results \& Empirical Evaluation}
% ------------------------------------------------
\begin{enumerate}
    \item Specify the benchmark dataset, data split, and resolution or setup used for training and testing.
    \item Present a comparative table contrasting your method against the reproduced baseline across standard evaluation metrics.
    \item State the primary empirical takeaway highlighting the exact performance delta achieved by your proposed contribution.
\end{enumerate}

% ------------------------------------------------
% Results Table
% ------------------------------------------------
\begin{center}

\captionof{table}{
Quantitative comparison of the baseline and
proposed method on [Dataset Name].
Best results are highlighted in bold.
}

\label{tab:results}

\small

\begin{tabularx}{\linewidth}{
    @{}X c c c c@{}
}

\toprule

\textbf{Method}
&
\makecell{\textbf{Metric 1}\\($\uparrow$)}
&
\makecell{\textbf{Metric 2}\\($\uparrow$)}
&
\makecell{\textbf{Params}\\(M)}
&
\makecell{\textbf{Latency}\\(ms $\downarrow$)}
\\

\midrule

Baseline
&
0.00
&
0.00
&
0.0
&
0.0
\\

\textbf{Ours}
&
\textbf{0.00}
&
\textbf{0.00}
&
0.0
&
0.0
\\

\midrule

$\Delta$ (Ours -- Baseline)
&
+0.00
&
+0.00
&
+0.0
&
+0.0
\\

\bottomrule

\end{tabularx}

\end{center}

\vspace{-5pt}

\textbf{Result Interpretation:}
Briefly explain the observed improvements
or performance trade-offs.

Explain how the proposed novelty contributed
to the experimental results and support
the interpretation with relevant evidence.

\end{document}
